% Этот текст можно свободно редактировать.
\FloatBarrier
\section{Сильный ручной baseline}

CatBoostClassifier, LogLoss как обучающая функция потерь. 48 признаков:

\begin{itemize}
\item Категории пользователя: пол, возрастная группа, доход, дети; объекта: ID, тип, основной жанр, страна, режиссёр, возрастной рейтинг и детский флаг. Год выпуска и день недели --- числовые. Пропуски выделяются явно.
\item Число прошлых записей пользователя и объекта за 7, 30 дней и всю историю.
\item Косвенные связи user--genre, user--country, user--director, user--content\_type: counts и доли в истории пользователя для тех же окон. Для этих основных агрегатов берётся первый токен соответствующей категории.
\item Давность последней записи пользователя и объекта, длина наблюдаемого периода пользователя, число различных просмотренных объектов, прямое число прошлых user--item взаимодействий.
\end{itemize}
Прямой \code{pair_count} при фильтрации уже виденных объектов всегда равен нулю. Он сохранён как явная проверка постановки; ожидать от него прироста здесь нельзя. Косвенные связи остаются содержательным источником персонализации.

Сглаженная доля реализована точно как \code{relation_count / (user_count + 1)}. Это описательный признак состава истории, без показов в знаменателе; его нельзя интерпретировать как CTR.

Проверены три набора \code{(depth,l2_leaf_reg)}: [\{'depth': 6, 'l2\_leaf\_reg': 3\}, \{'depth': 6, 'l2\_leaf\_reg': 10\}, \{'depth': 8, 'l2\_leaf\_reg': 10\}]. Выбран \textbf{\{'depth': 6, 'l2\_leaf\_reg': 3\}} по NDCG@10 на search. До 450 деревьев, learning rate 0.08, early stopping по AUC с patience 50. Используется CPU, 6 потоков. Все сравниваемые модели используют одинаковые гиперпараметры и доступную информацию. Пропорция контролей одинакова на train/search/confirm/test в основном режиме. User ID не используется как категориальный признак.

CatBoost обрабатывает категориальные признаки на train; данные валидации не используются для построения target statistics. Фактические параметры обученных моделей дополнительно проверяются: \code{counter_calc_method=SkipTest} исключает даже частотные Counter-статистики валидации; \code{boosting_type=Plain}, \code{one_hot_max_size=32}, \code{max_ctr_complexity=4}. Полные разрешённые параметры сохраняются в \code{catboost_resolved_parameters.json}. Обработка категорий и алгоритмы CatBoost описаны в \href{https://proceedings.neurips.cc/paper/2018/hash/14491b756b3a51daac41c24863285549-Abstract.html}{оригинальной статье}, а смысл SkipTest --- в \href{https://catboost.ai/docs/en/references/training-parameters/ctr}{официальной документации}. Важное допущение: демографические данные и метаданные трактуются как доступный статический справочник; исторических снимков этих таблиц нет. Формальная проверка временных агрегатов не устраняет эту неопределённость источника.

